\documentclass{article}
\input{../../../lib.sty}

\title{\texttt{fn truncate\_domain}}
\author{\MichaelShoemateAuthor}
\begin{document}
\maketitle


\contrib

Proves soundness of \rustdoc{transformations/make_stable_lazyframe/truncate/fn}{truncate\_domain}
in \asOfCommit{mod.rs}{a105afe}.

\section{Hoare Triple}
\subsection*{Precondition}
\subsubsection*{Compiler-verified}
None

\subsubsection*{Caller-verified}
Truncation keys are stable row-by-row transformations of the data.

\subsection*{Pseudocode}
\label{sec:python-pseudocode}
\lstinputlisting[language=Python,firstline=2,escapechar=|]{./pseudocode/truncate_domain.py}

\subsection*{Postcondition}
\begin{theorem}[Postcondition]
    Define \function as the function that takes in a dataset \texttt{arg} and returns a dataset with \texttt{truncation} applied to it.
    \texttt{truncation} is either a filter or groupby operation.
    \begin{itemize}
        \item If \texttt{truncation} is a filter, then \function\ returns a subset of rows in the data.
        \item If \texttt{truncation} is a groupby operation, then \function\ groups by \texttt{truncation.keys} and the identifier, and aggregates by \texttt{truncation.aggs}.
    \end{itemize}

    For every element $x$ in \texttt{domain}, $\function(x)$ is in the returned domain or raises a data-independent runtime exception.
\end{theorem}

\begin{proof}
    First consider the case where \texttt{truncation} is a filter.
    Since filtering is a contractive mapping,
    the output domain is a subset of the input domain.
    However, any invariants on the group length and keys are no longer valid.
    This is reflected in line \ref{line:invariant-filter}.

    Now consider the case where \texttt{truncation} is a groupby operation.
    Since each group becomes one row in the resulting dataset,
    an arbitrary black-box function may be applied to each group,
    so long as erors from the black-box function (the aggregates) are data-independent.
    Line \ref{line:infallible} ensures that the aggregates are infallible.

    The schema of members of the output domain is determined by the keys and aggregates of the groupby operation,
    which is computed on line \ref{line:simulate-schema}.
    Only margins whose grouping expressions are a subset of the original grouping keys are retained.
    For each such margin, line \ref{line:materialized-margin} replaces every key expression
    with \texttt{col(output\_name)} for that expression.
    A group key is evaluated on the input and stored under its output name by the groupby operation,
    so grouping output rows by these materialized columns uses the same labels as the original input grouping.
    Evaluating the original expression again on the output is not generally equivalent.

    Each original row contributes to at most one output row,
    and each output row retains its original grouping labels.
    Therefore the upper bounds on group lengths and the number of groups remain valid for the translated margins,
    even when the aggregates change other columns.
    Among preserved margins, invariants are discarded.
\end{proof}

\end{document}
